% =====================================================================
%  IE — Formulario de Instalaciones Electricas / Circuitos Magneticos
%  Compilar:  pdflatex IE.tex   (dos veces, por el indice)
%  Documento de respaldo teorico de la clase IE.m
% =====================================================================
\documentclass[11pt,a4paper]{article}

\usepackage[T1]{fontenc}
\usepackage[utf8]{inputenc}
\usepackage[spanish,es-tabla,es-noshorthands,es-nodecimaldot]{babel}
\usepackage[margin=2.3cm]{geometry}
\usepackage{amsmath,amssymb}
\usepackage{empheq}
\usepackage{siunitx}
\usepackage{booktabs}
\usepackage{array}
\usepackage{tikz}
\usepackage{enumitem}
\usepackage[colorlinks=true,linkcolor=black,urlcolor=blue]{hyperref}

\sisetup{per-mode=symbol, inter-unit-product=\ensuremath{{}\cdot{}}}

% unidades electricas que siunitx no trae de fabrica
\DeclareSIUnit{\voltampere}{VA}
\DeclareSIUnit{\var}{VAR}
\DeclareSIUnit{\vuelta}{vuelta}

% --- atajos ----------------------------------------------------------
\newcommand{\mur}{\mu_r}
\newcommand{\muz}{\mu_0}
\newcommand{\Rel}{\mathcal{R}}          % reluctancia
\newcommand{\fmm}{\mathcal{F}}          % fuerza magnetomotriz
\newcommand{\ju}{\mathrm{j}}
\newcommand{\fn}[1]{\texttt{#1}}        % nombre de funcion MATLAB
\newcommand{\Av}{\unit{\ampere}\text{-vuelta}}

% cabecera de tabla de variables
\newcolumntype{L}[1]{>{\raggedright\arraybackslash}p{#1}}

\title{\textbf{Formulario IE}\\[2mm]
\large Instalaciones Eléctricas y Circuitos Magnéticos\\
\normalsize Respaldo teórico de la caja de herramientas \fn{IE.m}}
\author{}
\date{}

\begin{document}
\maketitle
\vspace{-15mm}
\tableofcontents
\bigskip
\hrule
\bigskip

% =====================================================================
\section{Convenciones y constantes}
% =====================================================================

\begin{itemize}[nosep,leftmargin=*]
  \item Ángulos en \textbf{grados} en la interfaz; radianes solo dentro del cálculo.
  \item Tensiones y corrientes en valor \textbf{eficaz (RMS)} salvo aviso explícito.
  \item Secuencia de fases \textbf{positiva (ABC)}: B atrasa \ang{120} a A, C atrasa \ang{240}.
  \item $\mur$ es siempre la permeabilidad \textbf{relativa} (adimensional). El aire y el vacío tienen $\mur = 1$.
  \item Los modelos magnéticos son \textbf{lineales}: $\mur$ constante, sin curva $B$--$H$.
\end{itemize}

\begin{table}[h]
\centering
\begin{tabular}{@{}llll@{}}
\toprule
\textbf{Constante} & \textbf{Símbolo} & \textbf{Valor} & \textbf{Unidad} \\
\midrule
Permeabilidad del vacío   & $\muz$      & $4\pi\times10^{-7} \approx 1.2566\times10^{-6}$ & \unit{\henry\per\metre} \\
Permitividad del vacío    & $\varepsilon_0$ & $8.8542\times10^{-12}$ & \unit{\farad\per\metre} \\
Factor de cresta senoidal & $\sqrt{2}$  & $1.41421$ & --- \\
Factor trifásico          & $\sqrt{3}$  & $1.73205$ & --- \\
\bottomrule
\end{tabular}
\end{table}

\noindent\textbf{Permeabilidad absoluta:} $\mu = \mur\,\muz$ \; [\unit{\henry\per\metre}]. En todo el formulario
$\mur$ aparece siempre acompañada de $\muz$; nunca se usa la absoluta sola.

% =====================================================================
\section{Fasores, impedancias y potencia}
% =====================================================================

\subsection{Representación fasorial}

Una senoide $v(t) = V_{\max}\cos(\omega t + \theta)$ se representa por el fasor complejo RMS
\[
  \bar{V} = V \angle \theta = V\,e^{\ju\theta} = V\cos\theta + \ju\,V\sin\theta,
  \qquad V = \frac{V_{\max}}{\sqrt{2}}.
\]

\begin{align}
  \text{Polar} \to \text{Rectangular:}&\quad \bar{X} = x\,e^{\ju\theta} = x\cos\theta + \ju\,x\sin\theta \\
  \text{Rectangular} \to \text{Polar:}&\quad x = |\bar{X}| = \sqrt{a^2+b^2}, \qquad
                                             \theta = \arctan\!\left(\frac{b}{a}\right)
\end{align}

Relaciones entre valores de una senoide:
\[
  V_{\max} = \sqrt{2}\,V_{\text{RMS}}, \qquad
  V_{pp} = 2V_{\max} = 2\sqrt{2}\,V_{\text{RMS}}, \qquad
  \omega = 2\pi f .
\]

\begin{table}[h]
\centering
\begin{tabular}{@{}cL{7.6cm}c@{}}
\toprule
\textbf{Símbolo} & \textbf{Descripción} & \textbf{Unidad} \\
\midrule
$\bar{V},\bar{I}$ & Fasores de tensión y corriente (RMS, complejos) & \unit{\volt},\ \unit{\ampere} \\
$V_{\max}$        & Valor de pico (amplitud) & \unit{\volt} \\
$V_{pp}$          & Valor de pico a pico & \unit{\volt} \\
$\theta$          & Fase inicial & \unit{\degree} \\
$f$               & Frecuencia & \unit{\hertz} \\
$\omega$          & Pulsación, $\omega = 2\pi f$ & \unit{\radian\per\second} \\
\bottomrule
\end{tabular}
\end{table}

\subsection{Impedancia}

\[
  \bar{Z} = R + \ju X = |Z|\angle\varphi,
  \qquad
  X = X_L - X_C = \omega L - \frac{1}{\omega C},
  \qquad
  \bar{Y} = \frac{1}{\bar{Z}} = G + \ju B .
\]

\begin{table}[h]
\centering
\begin{tabular}{@{}cL{7.6cm}c@{}}
\toprule
\textbf{Símbolo} & \textbf{Descripción} & \textbf{Unidad} \\
\midrule
$\bar{Z}$ & Impedancia compleja & \unit{\ohm} \\
$R$       & Resistencia (parte real) & \unit{\ohm} \\
$X$       & Reactancia (parte imaginaria); $X>0$ inductiva, $X<0$ capacitiva & \unit{\ohm} \\
$\bar{Y}$ & Admitancia, inversa de $\bar{Z}$ & \unit{\siemens} \\
$L$       & Inductancia & \unit{\henry} \\
$C$       & Capacidad & \unit{\farad} \\
\bottomrule
\end{tabular}
\end{table}

\paragraph{Asociación.}
\begin{align}
  \text{Serie:}&\quad \bar{Z}_{eq} = \sum_{k=1}^{n} \bar{Z}_k \\
  \text{Paralelo:}&\quad \bar{Z}_{eq} = \left(\sum_{k=1}^{n} \frac{1}{\bar{Z}_k}\right)^{\!-1}
  \qquad\Longrightarrow\qquad
  \bar{Z}_{eq} = \frac{\bar{Z}_1\bar{Z}_2}{\bar{Z}_1+\bar{Z}_2} \ \ (n=2)
\end{align}

\paragraph{Divisores.} Con $\bar{V}$ aplicada a la serie y $\bar{I}$ entrando al nodo del paralelo:
\begin{align}
  \text{Tensión:}&\quad \bar{V}_k = \bar{V}\,\frac{\bar{Z}_k}{\displaystyle\sum_j \bar{Z}_j}
  &&\text{control:}\ \ \sum_k \bar{V}_k = \bar{V} \\[2mm]
  \text{Corriente:}&\quad \bar{I}_k = \bar{I}\,\frac{\bar{Y}_k}{\displaystyle\sum_j \bar{Y}_j}
                       = \bar{I}\,\frac{1/\bar{Z}_k}{\displaystyle\sum_j 1/\bar{Z}_j}
  &&\text{control:}\ \ \sum_k \bar{I}_k = \bar{I}
\end{align}

\subsection{Transformación estrella--triángulo (Kennelly)}

Es una transformación de \textbf{impedancias}, no de tensiones. El orden es
$\Delta = [\bar{Z}_{ab},\bar{Z}_{bc},\bar{Z}_{ca}]$ y $Y = [\bar{Z}_a,\bar{Z}_b,\bar{Z}_c]$,
donde $\bar{Z}_a$ es el brazo que cuelga del nodo $a$ hacia el neutro.

\paragraph{Caso general.} $\Delta \to Y$: cada brazo es el producto de las dos ramas que
tocan ese nodo, sobre la suma de las tres.
\begin{equation}
  \bar{Z}_a = \frac{\bar{Z}_{ab}\,\bar{Z}_{ca}}{\bar{Z}_{ab}+\bar{Z}_{bc}+\bar{Z}_{ca}},
  \qquad
  \bar{Z}_b = \frac{\bar{Z}_{ab}\,\bar{Z}_{bc}}{\bar{Z}_{ab}+\bar{Z}_{bc}+\bar{Z}_{ca}},
  \qquad
  \bar{Z}_c = \frac{\bar{Z}_{bc}\,\bar{Z}_{ca}}{\bar{Z}_{ab}+\bar{Z}_{bc}+\bar{Z}_{ca}}
\end{equation}

\noindent $Y \to \Delta$: cada rama es $P$ sobre el brazo \textbf{opuesto}, con
$P = \bar{Z}_a\bar{Z}_b + \bar{Z}_b\bar{Z}_c + \bar{Z}_c\bar{Z}_a$.
\begin{equation}
  \bar{Z}_{ab} = \frac{P}{\bar{Z}_c},
  \qquad
  \bar{Z}_{bc} = \frac{P}{\bar{Z}_a},
  \qquad
  \bar{Z}_{ca} = \frac{P}{\bar{Z}_b}
\end{equation}

\paragraph{Caso equilibrado.} Si las tres impedancias son iguales:
\begin{equation}
  \boxed{\ \bar{Z}_Y = \frac{\bar{Z}_\Delta}{3} \ }
  \qquad\qquad
  \boxed{\ \bar{Z}_\Delta = 3\,\bar{Z}_Y \ }
\end{equation}

\begin{quote}\small
\textbf{Cuidado con el 3 y el $\sqrt{3}$.} El factor $3$ es de \emph{impedancias};
el factor $\sqrt{3}$ es de \emph{tensiones} (sección siguiente). No son intercambiables.
La misma impedancia conectada en $\Delta$ consume tres veces la potencia que en $Y$:
por eso al pasar a $Y$ hay que dividir por 3 para conservar el comportamiento.
\end{quote}

\subsection{Sistema trifásico equilibrado}

Con secuencia positiva y $\bar{E}_{AN} = E\angle\theta_0$:
\begin{equation}
  \bar{E}_{AN} = E\angle\theta_0, \qquad
  \bar{E}_{BN} = E\angle(\theta_0 - \ang{120}), \qquad
  \bar{E}_{CN} = E\angle(\theta_0 - \ang{240}) .
\end{equation}

Las tensiones de línea salen de la diferencia de dos de fase:
\begin{equation}
  \bar{E}_{AB} = \bar{E}_{AN} - \bar{E}_{BN} = \sqrt{3}\,E \angle(\theta_0 + \ang{30}) .
\end{equation}

\noindent Es decir, la línea es $\sqrt{3}$ veces mayor \emph{y} adelanta \ang{30} a su fase homóloga:
\begin{equation}
  \boxed{\ E_L = \sqrt{3}\,E_F \ }
  \qquad\qquad
  \boxed{\ E_F = \frac{E_L}{\sqrt{3}} \ }
\end{equation}

\begin{table}[h]
\centering
\begin{tabular}{@{}lccc@{}}
\toprule
\textbf{Conexión} & \textbf{Tensión} & \textbf{Corriente} & \textbf{Potencia activa} \\
\midrule
Estrella ($Y$)     & $E_L = \sqrt{3}\,E_F$ & $I_L = I_F$            & $P = \sqrt{3}\,E_L I_L\cos\varphi$ \\
Triángulo ($\Delta$) & $E_L = E_F$          & $I_L = \sqrt{3}\,I_F$ & $P = \sqrt{3}\,E_L I_L\cos\varphi$ \\
\bottomrule
\end{tabular}
\end{table}

\noindent En estrella la corriente de línea \emph{es} la de fase; en triángulo lo que coincide
es la tensión. La expresión de la potencia en función de magnitudes de línea es la misma
para las dos conexiones: por eso conviene trabajar siempre en variables de línea.
La potencia total es $3$ veces la de una fase: $P_{3\phi} = 3\,E_F I_F \cos\varphi$.

\subsection{Potencia}

Con fasores RMS, la potencia compleja es
\begin{equation}
  \boxed{\ \bar{S} = \bar{V}\,\bar{I}^{*} = P + \ju Q \ }
\end{equation}

\begin{align}
  P     &= \Re\{\bar{S}\} = V I \cos\varphi &&\text{activa}      & [\unit{\watt}] \\
  Q     &= \Im\{\bar{S}\} = V I \sin\varphi &&\text{reactiva}    & [\unit{\var}] \\
  |S|   &= |\bar{V}||\bar{I}| = \sqrt{P^2+Q^2} &&\text{aparente} & [\unit{\voltampere}] \\
  \mathrm{FP} &= \cos\varphi = \frac{P}{|S|} &&\text{factor de potencia} & [-]
\end{align}

\noindent El ángulo del triángulo es $\varphi = \arg(\bar{S}) = \theta_V - \theta_I$.

\begin{table}[h]
\centering
\begin{tabular}{@{}lccl@{}}
\toprule
\textbf{Signo de $Q$} & \textbf{Carácter} & \textbf{Corriente} & \textbf{Ejemplo} \\
\midrule
$Q > 0$ & Inductiva  & atrasa a la tensión   & motor, bobina \\
$Q = 0$ & Resistiva  & en fase               & estufa, lámpara incandescente \\
$Q < 0$ & Capacitiva & adelanta a la tensión & banco de capacitores \\
\bottomrule
\end{tabular}
\end{table}

\paragraph{Por qué el conjugado.} Si se escribiera $\bar{S} = \bar{V}\bar{I}$ el ángulo
resultante sería $\theta_V + \theta_I$, que no tiene significado físico. El conjugado hace
aparecer la \emph{diferencia} $\theta_V - \theta_I$, que es el desfase real entre tensión y
corriente y por lo tanto el ángulo del triángulo de potencias.

\paragraph{Corrección del factor de potencia.} Para llevar el FP de $\cos\varphi_1$ a
$\cos\varphi_2$ manteniendo $P$:
\begin{equation}
  Q_C = P\,(\tan\varphi_1 - \tan\varphi_2),
  \qquad
  C = \frac{Q_C}{\omega V^2} \quad [\unit{\farad}].
\end{equation}

% =====================================================================
\section{Circuito magnético}
% =====================================================================

\subsection{Variables y unidades}

\begin{table}[h]
\centering
\begin{tabular}{@{}cL{7.4cm}ll@{}}
\toprule
\textbf{Símbolo} & \textbf{Descripción} & \textbf{Unidad} & \textbf{Equivalencia} \\
\midrule
$B$    & Densidad de flujo magnético (inducción) & \unit{\tesla}   & \unit{\weber\per\metre\squared} \\
$H$    & Intensidad de campo magnético           & \unit{\ampere\per\metre} & \Av/\unit{\metre} \\
$\Phi$ & Flujo magnético                         & \unit{\weber}   & \unit{\volt\second} \\
$\dot{\Phi}$ & Derivada del flujo, $\mathrm{d}\Phi/\mathrm{d}t$ & \unit{\weber\per\second} & \unit{\volt} \\
$\mur$ & Permeabilidad relativa del material     & ---             & $\mu/\muz$ \\
$\muz$ & Permeabilidad del vacío                 & \unit{\henry\per\metre} & $4\pi\times10^{-7}$ \\
$\ell$ & Longitud media del camino magnético     & \unit{\metre}   & $2\pi R$ en un toroide \\
$A$    & Sección transversal del núcleo          & \unit{\metre\squared} & $a\cdot b$ si es rectangular \\
$N$    & Número de vueltas de la bobina          & ---             & \\
$I$    & Corriente por la bobina                 & \unit{\ampere}  & \\
$\fmm$ & Fuerza magnetomotriz, $NI$              & \Av             & \\
$\Rel$ & Reluctancia                             & \Av/\unit{\weber} & \unit{\per\henry} \\
$\mathcal{P}$ & Permeancia, $1/\Rel$             & \unit{\henry}   & \\
$g$    & Longitud del entrehierro                & \unit{\metre}   & \\
$\delta_A$ & Crecimiento proporcional del área en el gap & ---     & \\
$L$    & Inductancia                             & \unit{\henry}   & \unit{\weber\per\ampere} \\
$F$    & Fuerza de atracción en el entrehierro   & \unit{\newton}  & \\
$v$    & Tensión inducida (f.e.m.)               & \unit{\volt}    & \\
\bottomrule
\end{tabular}
\end{table}

\subsection{Las cuatro leyes básicas}

Éste es el sistema que resuelve \fn{IE.datosB}:

\begin{empheq}[box=\fbox]{align}
  H  &= \frac{N I}{\ell}      &&\text{Ley de Ampère} \label{eq:ampere}\\
  B  &= \mur\,\muz\,H         &&\text{Relación } B\text{--}H \text{ del material (lineal)} \label{eq:bh}\\
  B  &= \frac{\Phi}{A}        &&\text{Definición de densidad de flujo} \label{eq:phi}\\
  v  &= N\,\frac{\mathrm{d}\Phi}{\mathrm{d}t} &&\text{Ley de Faraday} \label{eq:faraday}
\end{empheq}

\noindent Combinando \eqref{eq:ampere}--\eqref{eq:phi}:
\begin{equation}
  \Phi = \frac{\mur\,\muz\,N I A}{\ell}
       = \frac{NI}{\ell/(\mur\muz A)}
       = \frac{\fmm}{\Rel}.
\end{equation}

\begin{quote}\small
\textbf{Por qué Ampère y $B$--$H$ van separadas.} Fusionarlas en
$B = \mur\muz NI/\ell$ obliga a que $\mur$ cargue dos significados distintos
(relativa aquí, absoluta en $\mu = B/H$) y el sistema queda contradictorio.
Dejarlas separadas mantiene una sola definición de $\mur$.
\end{quote}

\subsection{Ley de Hopkinson y reluctancia}

La analogía con el circuito eléctrico es exacta mientras el modelo sea lineal:

\begin{equation}
  \boxed{\ \fmm = \Phi\,\Rel \ }
  \qquad\text{(análoga a } V = I R\text{)}
  \qquad\qquad
  \Rel = \frac{\ell}{\mur\,\muz\,A}
\end{equation}

\begin{table}[h]
\centering
\begin{tabular}{@{}lll@{}}
\toprule
& \textbf{Circuito eléctrico} & \textbf{Circuito magnético} \\
\midrule
Excitación   & Tensión $V$ \hfill [\unit{\volt}]      & F.m.m. $\fmm = NI$ \hfill [\Av] \\
Respuesta    & Corriente $I$ \hfill [\unit{\ampere}]  & Flujo $\Phi$ \hfill [\unit{\weber}] \\
Oposición    & Resistencia $R = \rho\ell/A$ \hfill [\unit{\ohm}] & Reluctancia $\Rel = \ell/(\mur\muz A)$ \hfill [\Av/\unit{\weber}] \\
Inversa      & Conductancia $G$ \hfill [\unit{\siemens}] & Permeancia $\mathcal{P}$ \hfill [\unit{\henry}] \\
Ley de mallas& $\sum V = 0$                           & $\sum \fmm = 0$ (Ampère) \\
Ley de nodos & $\sum I = 0$                           & $\sum \Phi = 0$ \\
En serie     & $R_{eq} = \sum R_k$                    & $\Rel_{eq} = \sum \Rel_k$ \\
En paralelo  & $1/R_{eq} = \sum 1/R_k$                & $1/\Rel_{eq} = \sum 1/\Rel_k$ \\
\bottomrule
\end{tabular}
\end{table}

\noindent \textbf{Dónde se rompe la analogía:} $\rho$ es constante, $\mur$ \emph{no}.
Al saturar, $\mur$ se desploma y $\Rel$ crece. La analogía vale solo por debajo del codo
de la curva $B$--$H$.

\subsection{Entrehierro y dispersión de bordes (\emph{fringing})}

Un tramo con entrehierro se descompone en dos reluctancias en serie: hierro y aire.
La longitud del gap se \textbf{descuenta} de la del hierro.

\begin{equation}
  \Rel_n = \frac{\ell - g}{\mur\,\muz\,A},
  \qquad
  \Rel_g = \frac{g}{\muz\,A_g},
  \qquad
  \Rel_s = \Rel_n + \Rel_g .
\end{equation}

Al cruzar el entrehierro el flujo se abre, así que la sección efectiva del aire es
\textbf{mayor} que la del hierro:
\begin{equation}
  \boxed{\ A_g = A\,(1 + \delta_A) \ }
\end{equation}

\noindent Si se conoce la geometría rectangular $a \times b$, la regla estándar es agregar
$g$ a cada lado:
\begin{equation}
  A_g = (a+g)(b+g)
  \qquad\Longrightarrow\qquad
  \delta_A = \frac{(a+g)(b+g)}{a\,b} - 1 .
\end{equation}

Como el flujo es el mismo pero las secciones difieren, las densidades también:
\begin{equation}
  B = \frac{\Phi}{A} \quad\text{(hierro)},
  \qquad\qquad
  B_g = \frac{\Phi}{A_g} = \frac{B}{1+\delta_A} \quad\text{(entrehierro)} .
\end{equation}

\noindent Es decir $B_g < B$: el entrehierro tiene más área, así que menos densidad.

\begin{quote}\small
\textbf{El entrehierro domina.} Con $\mur \sim 5000$, un gap de \qty{1}{\milli\metre} en un
camino de \qty{300}{\milli\metre} aporta más del \qty{90}{\percent} de la reluctancia total.
Un milímetro de aire equivale a varios metros de hierro. Por eso el gap es la variable de
diseño que gobierna el circuito.
\end{quote}

\subsection{Inductancia y ley de Faraday}

\begin{equation}
  \lambda = N\Phi \quad [\unit{\weber}\text{-vuelta}],
  \qquad
  \boxed{\ L = \frac{\lambda}{I} = \frac{N\Phi}{I} = \frac{N^2}{\Rel} = \frac{\mur\muz N^2 A}{\ell} \ }
\end{equation}

\noindent La inductancia crece con $N^2$: una es por el flujo que la bobina genera, la otra
por las vueltas que ese flujo concatena.

\paragraph{Tensión inducida.} Las dos formas son equivalentes:
\begin{equation}
  v(t) = N\,\frac{\mathrm{d}\Phi}{\mathrm{d}t}
       = L\,\frac{\mathrm{d}i}{\mathrm{d}t} .
\end{equation}

\paragraph{Caso senoidal.} Con $\Phi(t) = \Phi_{\max}\sin(\omega t)$:
\begin{equation}
  v(t) = N\omega\Phi_{\max}\cos(\omega t)
  \qquad\Longrightarrow\qquad
  V_{\text{RMS}} = \frac{N\omega\Phi_{\max}}{\sqrt{2}}
                 = \boxed{\,4.44\,N f \Phi_{\max}\,}
\end{equation}
donde $4.44 = 2\pi/\sqrt{2}$. Es la \textbf{fórmula fundamental del transformador}.

\paragraph{Energía almacenada.}
\begin{equation}
  W = \tfrac{1}{2} L I^2 \quad [\unit{\joule}],
  \qquad
  w = \frac{B^2}{2\muz} \quad [\unit{\joule\per\metre\cubed}] \ \text{(densidad en el aire)} .
\end{equation}

\subsection{Fuerza en el entrehierro}

\begin{equation}
  \boxed{\ F = \frac{B_g^2\,A_g}{2\,\muz} \ } \qquad [\unit{\newton}]
\end{equation}

\noindent \textbf{Por cada superficie de contacto.} Un electroimán en U con una barra
encima tiene \emph{dos} contactos: la fuerza total es el doble. La fuerza va con $B_g^2$,
así que duplicar la corriente cuadruplica la fuerza (mientras no sature).

\subsection{Núcleo tipo O — camino cerrado en serie}

\begin{center}
\begin{tikzpicture}[scale=0.85]
  \draw[line width=6pt] (0,0) rectangle (4,2.6);
  \node at (2,-0.6) {$\Rel_t = \Rel_1+\Rel_2+\Rel_3+\Rel_4 \ (+\,\Rel_g)$};
  \node[above] at (2,1.4) {$\Phi$};
  \draw[->,line width=1pt] (1.3,1.3) -- (2.7,1.3);
  \node[left] at (-0.35,1.3) {$N,\,I$};
  \foreach \y in {0.7,1.0,1.3,1.6,1.9} \draw[line width=1pt] (-0.28,\y) arc (90:270:0.13);
\end{tikzpicture}
\end{center}

Cuadrado, rectangular, toroide, U+I, E-I armado: todos son un solo lazo. El flujo no se
divide y las reluctancias se suman.

\begin{align}
  \fmm &= N I  & \Rel_t &= \sum_k \Rel_k \\
  \Phi &= \frac{\fmm}{\Rel_t} = \frac{NI}{\sum_k \Rel_k}
       & L &= \frac{N^2}{\Rel_t}
\end{align}

\noindent Caída de f.m.m. en cada tramo y control de Ampère:
\begin{equation}
  \fmm_k = \Phi\,\Rel_k
  \qquad\qquad
  \boxed{\ \sum_k \fmm_k = \fmm = NI \ }
\end{equation}

\noindent Ese control es el análogo magnético de la ley de mallas y es la verificación más
sensible del cálculo: si no cierra, hay un error en las reluctancias.

\subsection{Núcleo tipo E — dos caminos en paralelo}

\begin{center}
\begin{tikzpicture}[scale=0.85]
  \draw[line width=6pt] (0,0) rectangle (5,2.6);
  \draw[line width=6pt] (2.5,0) -- (2.5,2.6);
  \node at (2.5,-1.0) {$\Rel_t = \Rel_b + \left(\Rel_{o_1} \parallel \Rel_{o_2}\right)$};
  \node[above] at (2.5,2.75) {yugo};
  \foreach \y in {0.7,1.0,1.3,1.6,1.9} \draw[line width=1pt] (2.28,\y) arc (90:270:0.13);
  \node[left] at (2.05,1.3) {$N,\,I$};
  \node[below] at (0,-0.15) {$o_1$};
  \node[below] at (2.5,-0.15) {$b$};
  \node[below] at (5,-0.15) {$o_2$};
\end{tikzpicture}
\end{center}

La bobina va en una pierna $b$; el flujo vuelve repartido por las otras dos, que quedan en
paralelo.

\begin{align}
  \Rel_t &= \Rel_b + \left(\frac{1}{\Rel_{o_1}} + \frac{1}{\Rel_{o_2}}\right)^{\!-1}
  & \Phi_b &= \frac{N I}{\Rel_t} \\[2mm]
  \Phi_{o_k} &= -\,\Phi_b\,\frac{1/\Rel_{o_k}}{1/\Rel_{o_1}+1/\Rel_{o_2}}
  & \sum_k \Phi_k &= 0
\end{align}

\noindent El reparto es el mismo divisor que el de corriente: inverso a la reluctancia.
El signo negativo indica que el flujo \emph{vuelve} por esas piernas; midiendo todos los
flujos entrando al yugo superior, la suma es cero (ley de nodos magnética).

\paragraph{Qué desprecia este modelo.}
\begin{enumerate}[nosep,leftmargin=*]
  \item Los \textbf{yugos} (barras horizontales) se suponen de reluctancia despreciable.
        Vale si son cortos y gruesos; deja de valer si hay entrehierro en el yugo.
  \item \textbf{Flujo de dispersión} entre piernas: cero. Todo lo que sale de $b$ vuelve
        por las otras dos.
  \item Es \textbf{lineal}: sin saturación.
\end{enumerate}

% =====================================================================
\section{Transformador ideal}
% =====================================================================

\subsection{Origen: Faraday sobre un núcleo común}

Las dos bobinas están arrolladas sobre el mismo núcleo, así que ven el
\textbf{mismo} $\mathrm{d}\Phi/\mathrm{d}t$:
\begin{equation}
  V_p = N_p\,\frac{\mathrm{d}\Phi}{\mathrm{d}t},
  \qquad
  V_s = N_s\,\frac{\mathrm{d}\Phi}{\mathrm{d}t}
  \qquad\Longrightarrow\qquad
  \frac{V_p}{V_s} = \frac{N_p}{N_s} .
\end{equation}

En el núcleo ideal $\Rel \to 0$, con lo cual la f.m.m. neta debe ser nula:
\begin{equation}
  N_p I_p - N_s I_s = \Phi\,\Rel = 0
  \qquad\Longrightarrow\qquad
  \frac{I_s}{I_p} = \frac{N_p}{N_s} .
\end{equation}

\subsection{Relación de transformación}

\begin{equation}
  \boxed{\ a = \frac{N_p}{N_s} \ } \quad\text{(constante primaria)}
  \qquad\qquad
  \boxed{\ b = \frac{N_s}{N_p} = \frac{1}{a} \ } \quad\text{(constante secundaria)}
\end{equation}

\begin{table}[h]
\centering
\begin{tabular}{@{}ccc@{}}
\toprule
\textbf{Magnitud} & \textbf{Referida al primario} & \textbf{Referida al secundario} \\
\midrule
Tensión    & $V_p = a\,V_s$    & $V_s = b\,V_p$ \\
Corriente  & $I_p = I_s/a$     & $I_s = I_p/b$ \\
Impedancia & $Z_p = a^2\,Z_s$  & $Z_s = b^2\,Z_p$ \\
\bottomrule
\end{tabular}
\end{table}

\noindent \textbf{La tensión y la corriente van con $a$; la impedancia con $a$ al
cuadrado.} Ese cuadrado es el que más se olvida. Sale de dividir las dos primeras:
\begin{equation}
  Z_p = \frac{V_p}{I_p} = \frac{a\,V_s}{I_s/a} = a^2\,\frac{V_s}{I_s} = a^2 Z_s .
\end{equation}

\begin{table}[h]
\centering
\begin{tabular}{@{}cL{7.2cm}c@{}}
\toprule
\textbf{Símbolo} & \textbf{Descripción} & \textbf{Unidad} \\
\midrule
$a$      & Relación de transformación, $N_p/N_s$ & --- \\
$b$      & Relación inversa, $N_s/N_p = 1/a$ & --- \\
$N_p,N_s$& Vueltas de primario y secundario & --- \\
$V_p,V_s$& Tensiones de primario y secundario & \unit{\volt} \\
$I_p,I_s$& Corrientes de primario y secundario & \unit{\ampere} \\
$Z_p$    & Impedancia \emph{vista} desde el primario & \unit{\ohm} \\
$Z_s$    & Impedancia de carga conectada al secundario & \unit{\ohm} \\
$\dot{\Phi}$ & $\mathrm{d}\Phi/\mathrm{d}t$, común a las dos bobinas & \unit{\weber\per\second} \\
\bottomrule
\end{tabular}
\end{table}

\noindent Criterio rápido: $a > 1$ es \textbf{reductor} (baja la tensión, sube la
corriente); $a < 1$ es \textbf{elevador}.

\subsection{Conservación de la potencia}

En el transformador ideal no hay pérdidas ni almacenamiento, así que entra lo mismo que sale:
\begin{equation}
  \bar{S}_p = \bar{V}_p \bar{I}_p^{*}
            = (a\bar{V}_s)\left(\frac{\bar{I}_s}{a}\right)^{\!*}
            = \bar{V}_s \bar{I}_s^{*} = \bar{S}_s .
\end{equation}
Como $a$ es real, referir no rota fases: el factor de potencia visto desde el primario es
el mismo que el de la carga.

\subsection{Dimensionamiento del núcleo}

De la fórmula fundamental de la sección 3.5, con $\Phi_{\max} = B_{\max} A$:
\begin{equation}
  V = 4.44\,N f\,B_{\max}\,A
  \qquad\Longrightarrow\qquad
  N_p = \frac{V_p}{4.44\,f\,B_{\max}\,A} .
\end{equation}
$B_{\max}$ se elige por debajo de la saturación (típicamente \qtyrange{1.2}{1.6}{\tesla} en
chapa de hierro-silicio).

\subsection{Qué significa ``ideal''}

\begin{enumerate}[nosep,leftmargin=*]
  \item Sin resistencia de bobinados ($R_p = R_s = 0$): no hay pérdidas en el cobre.
  \item Sin flujo de dispersión ($X_p = X_s = 0$): todo el flujo concatena ambas bobinas.
  \item Sin pérdidas en el hierro (histéresis ni corrientes de Foucault).
  \item Corriente de magnetización nula ($\mur \to \infty$, $\Rel \to 0$).
\end{enumerate}

\noindent Consecuencia: rendimiento \qty{100}{\percent} y regulación de tensión nula. Si el
problema pide \emph{regulación} o \emph{rendimiento}, este modelo no alcanza — hay que
agregar la rama de excitación y las impedancias serie $R_{eq}$, $X_{eq}$ del circuito
equivalente real.

% =====================================================================
\section{Límites de los modelos}
% =====================================================================

\begin{enumerate}[leftmargin=*]
  \item \textbf{Todo el modelo magnético es lineal.} $\mur$ es constante y no hay curva
        $B$--$H$. Si algún tramo supera $\sim\qty{1.5}{\tesla}$, el resultado es
        \emph{optimista}: en la realidad $\mur$ se desploma al saturar y el flujo real es
        \textbf{menor} que el calculado. Con saturación hay que ir a la curva $B$--$H$.
  \item \textbf{El núcleo E ignora los yugos} y la dispersión entre piernas.
  \item \textbf{$E_L = \sqrt{3}E_F$ es solo de módulos}; el desfase de \ang{30} hay que
        aplicarlo aparte.
  \item \textbf{El sistema trifásico se asume equilibrado} y en secuencia positiva.
  \item \textbf{La potencia calculada es por fase.} Para la trifásica total hay que
        multiplicar por 3.
  \item \textbf{El transformador es ideal}: sin pérdidas, sin dispersión, sin corriente de
        magnetización.
\end{enumerate}

% =====================================================================
\appendix
\section{Correspondencia con \fn{IE.m}}
% =====================================================================

\begin{table}[h]
\centering
\small
\begin{tabular}{@{}lL{6.4cm}l@{}}
\toprule
\textbf{Función} & \textbf{Qué calcula} & \textbf{Sección} \\
\midrule
\fn{pol2rec}     & $x\angle\theta \to a + \ju b$ & 2.1 \\
\fn{parZ}        & Impedancias en paralelo & 2.2 \\
\fn{divV}, \fn{divI} & Divisores de tensión y corriente & 2.2 \\
\fn{d2e}, \fn{e2d}   & Transformación $\Delta \leftrightarrow Y$ (Kennelly) & 2.3 \\
\fn{f2l}, \fn{l2f}   & $E_L = \sqrt{3}E_F$ y su inversa (módulos) & 2.4 \\
\fn{fasores}     & Tabla completa del trifásico equilibrado & 2.4 \\
\fn{potencias}   & Triángulo de potencias, $\bar{S} = \bar{V}\bar{I}^{*}$ & 2.5 \\
\midrule
\fn{ecuacionesB}, \fn{datosB} & Las cuatro leyes básicas & 3.2 \\
\fn{despejar}    & Motor de sustitución hacia adelante & --- \\
\fn{ecuacionesSeg}, \fn{datosSeg} & Tramo con entrehierro & 3.4 \\
\fn{reluctSeg}   & Reluctancia de un tramo, con \emph{fringing} & 3.4 \\
\fn{faradayB}    & Análisis temporal: $v(t)$, $\Phi(t)$, $L$ & 3.5 \\
\fn{fluxO}       & Núcleo cerrado en serie, con fuerza y saturación & 3.7 \\
\fn{fluxE}       & Núcleo de tres piernas & 3.8 \\
\midrule
\fn{ecuacionesTrafo}, \fn{datosTrafo} & Modelo simbólico del transformador ideal & 4 \\
\fn{trafo}       & Resuelve todo el transformador de una & 4 \\
\bottomrule
\end{tabular}
\end{table}

\vfill
\noindent\hrulefill\\
\small Documento de respaldo de \fn{IE.m}. Las convenciones de signos, unidades y
secuencia de fases son las mismas en ambos.

\end{document}
